\subsection{例2：自由模的情形}

设 \(k\) 为一个环， \(\mathbf{M}\) 为 \(k\) -模， \(\mathbf{I}\) 为一个集合， \(p\) 为一个整数 \(\geqslant 0\) 。称映射 \(\boldsymbol{m}:1^{p}\to\mathbf{M}\) 是交错的，如果它满足以下两个条件：

\begin{enumerate}
\def\labelenumi{\alph{enumi})}
\tightlist
\item
  对每个置换 \(\sigma \in \mathfrak{S}_p\) 以及每个序列 \((\alpha_1, \ldots, \alpha_p) \in \mathrm{I}^p\) ，我们有
\end{enumerate}

\[
\boldsymbol {m} \left(\alpha_ {\sigma (1)}, \dots , \alpha_ {\sigma (p)}\right) = \varepsilon_ {\sigma} \boldsymbol {m} \left(\alpha_ {1}, \dots , \alpha_ {p}\right),
\]

\begin{enumerate}
\def\labelenumi{\alph{enumi})}
\setcounter{enumi}{1}
\tightlist
\item
  对每个序列 \((\alpha_{1},\ldots,\alpha_{p})\in\mathbb{I}^{p}\) ，若指标 \(\alpha_{1},\ldots,\alpha_{p}\) 中有两个相等，则 \(m(\alpha_{1},\ldots,\alpha_{p})=0\) 。
\end{enumerate}

（在 I 是 k-模且 m 是多线性的情形下，就恢复了 III，第 80 页中引入的概念。）

假设 I 是有限的，并用 \(C_{I}^{p}(M)\) 表示从 \(I^{p}\) 到 M 的交错映射的 k-模。

设 \(L_{0}\) 是一个 k-模， \((e_{i})_{i\in I}\) 为属于 \(L_{0}\) 的一个元素族；我们定义两个 k-线性映射

\[
\begin{array}{l} g: \operatorname{Hom} _ {k} (\symbf {\Lambda} ^ {p} \mathrm{L} _ {0}, \mathrm{M}) \to \mathrm{C} _ {\mathrm{I}} ^ {p} (\mathrm{M}) \\ h: \mathrm{C} _ {\mathrm{I}} ^ {p} (\mathrm{M}) \to \mathrm{M} \otimes_ {k} \symbf {\Lambda} ^ {p} \mathrm{L} _ {0} \end{array}
\]

如下：如果 \(f \in \operatorname{Hom}_k(\symbf{\Lambda}^p \mathbf{L}_0, \mathbf{M})\) ，我们设

\[
g (f) (\alpha_ {1}, \dots , \alpha_ {p}) = f (e _ {\alpha_ {1}} \wedge \dots \wedge e _ {\alpha_ {p}});
\]

设 \(m \in C_{I}^{p}(M)\) ，我们定义 \(h(m) \in M \otimes_{k} \Lambda^{p} L_{0}\) 。对每个元素 \((\alpha_{1}, \ldots, \alpha_{p})\)

属于 \(I^{p}\) ，元素 \(m(\alpha_{1},\ldots,\alpha_{p})\otimes(e_{\alpha_{1}}\wedge\ldots\wedge e_{\alpha_{p}})\) （属于 \(\Lambda^{p}L_{0}\otimes_{k}M\) ）当 Card \(\{\alpha_{1},\ldots,\alpha_{p}\}<p\) 时为零，并且与指标 \(\alpha_{1},\ldots,\alpha_{p}\) 的顺序无关，如果

\[
\text { Card } \{\alpha_ {1}, \dots , \alpha_ {p} \} = p.
\]

的顺序无关。它只依赖于 I 的子集 \(\mathbf{J} = (\alpha_{1},\ldots ,\alpha_{p})\) ；将其记作 \(h_\mathrm{J}(\pmb {m})\) ；我们有 \(h_\mathrm{J}(\pmb {m}) = 0\) ，如果 Card (J) \(< p\) ；然后我们设：

\[
h (\boldsymbol {m}) = \sum_ {\mathrm{J}} h _ {\mathrm{J}} (\boldsymbol {m}),
\]

其中J遍历I的所有含p个元素的子集。

引理 2. --- 若 k-模 \(\mathbf{L}_{0}\) 以 \((e_{i})_{i\in I}\) 为基，则 \(k\) -线性映射 \(g\) 与 \(h\) 都是双射的。

这由 III，第 79 页，定理 1 得出。

现在设 M 为一个 k-模，且 \(\boldsymbol{x} = (x_{i})_{i \in I}\) 为 M 的一族两两交换的 k-自同态。考虑多项式环 \(\mathbf{A} = k[(\mathbf{X}_{i})_{i \in I}]\) 并赋予 M 以 A-模结构，使得 \(\mathrm{P}(\mathbf{X}_{i}) m = \mathrm{P}(x_{i}) m\) 对 \(P \in A\) 和 \(m \in M\) 成立。另一方面，设 L 为自由 A-模 \(A^{I}, (e_{i})_{i \in I}\) 为其典范基， \(u : L \to A\) 为把 \(e_{i}\) 映到 \(X_{i}\) （对所有 \(i \in I\) ）的线性形式。考虑 k-模复形 \(\mathbf{K}^{\mathbf{A}}(u, \mathbf{M})\) 与 \(\mathbf{K}_{\mathbf{A}}^{\bullet}(u, \mathbf{M})\) ；我们有典范同构

\[
\begin{array}{c} \mathbf {K} _ {\mathbf {A}} ^ {p} (u, \mathrm{M}) = \operatorname{Hom} _ {\mathbf {A}} \left(\boldsymbol {\Lambda} _ {\mathbf {A}} ^ {p} (\mathrm{A} ^ {\mathrm{I}}), \mathrm{M}\right) \to \operatorname{Hom} _ {k} \left(\boldsymbol {\Lambda} _ {k} ^ {p} (k ^ {\mathrm{I}}), \mathrm{M}\right), \\ \mathrm{M} \otimes_ {k} \boldsymbol {\Lambda} _ {k} ^ {p} (k ^ {\mathrm{I}}) \to \mathrm{M} \otimes_ {\mathbf {A}} \boldsymbol {\Lambda} _ {\mathbf {A}} ^ {p} (\mathrm{A} ^ {\mathrm{I}}) = \mathbf {K} _ {p} ^ {\mathbf {A}} (u, \mathrm{M}); \end{array}
\]

由此，通过与同构 \(g\) 和 \(h\) 复合，得到 \(k\) -模的同构

\[
\begin{array}{l} \theta^ {p}: \mathbf {K} _ {\mathrm{A}} ^ {p} (u, \mathrm{M}) \to \mathrm{C} _ {\mathrm{I}} ^ {p} (\mathrm{M}), \\ \theta_ {p}: \mathrm{C} _ {\mathrm{I}} ^ {p} (\mathrm{M}) \to \mathbf {K} _ {p} ^ {\mathrm{A}} (u, \mathrm{M}). \end{array}
\]

我们用

\[
\begin{array}{l} \hat {c} ^ {p}: \mathbf {C} _ {\mathrm{I}} ^ {p} (\mathbf {M}) \to \mathbf {C} _ {\mathrm{I}} ^ {p + 1} (\mathbf {M}), \\ \hat {c} _ {p}: \mathbf {C} _ {\mathrm{I}} ^ {p} (\mathbf {M}) \to \mathbf {C} _ {\mathrm{I}} ^ {p - 1} (\mathbf {M}), \end{array}
\]

的 \(k\) -同态，它们是通过搬运 \(\mathbf{K}_{\mathbf{A}}^{\bullet}(u, \mathbf{M})\) 与 \(\mathbf{K}_{\mathbf{A}}^{\mathbf{A}}(u, \mathbf{M})\) 的微分经同构 \(\theta\) 而得到的。例如，我们有：

\[
\left(\partial^ {p} \boldsymbol {m}\right) \left(\alpha_ {1}, \dots , \alpha_ {p + 1}\right) = \sum_ {j = 1} ^ {p + 1} (- 1) ^ {j + 1} x _ {\alpha_ {j}} \boldsymbol {m} \left(\alpha_ {1}, \dots , \alpha_ {j - 1}, \alpha_ {j + 1}, \dots , \alpha_ {p + 1}\right).\tag{12}
\]

由 \(C_{I}^{p}(M)\) 与 \(\partial^{p}\) （相应地， \(\partial_{p}\) ）构成的复形记作 \(\mathbf{K}^{\prime}(\mathbf{x},\mathbf{M})\) （相应地， \(\mathbf{K}(\mathbf{x},\mathbf{M})\) ），称为与模 M 及自同态序列

\((x_{1}, \ldots, x_{n})\) 相伴的上升（相应地，下降）科斯居尔（Koszul）复形。因此我们有 k-模复形的同构

\[
\begin{array}{l} \theta^ {\bullet}: \mathbf {K} _ {\mathbf {A}} ^ {\bullet} (u, \mathbf {M}) \to \mathbf {K} ^ {\bullet} (x, \mathbf {M}), \\ \theta_ {\bullet}: \mathbf {K}. (x, \mathbf {M}) \to \mathbf {K} _ {\bullet} ^ {\mathbf {A}} (u, \mathbf {M}). \end{array}
\]

注记。--- 反过来，设 B 是一个 \(k\) -代数，L 是一个以 \((e_i)_{i \in I}\) 为基的自由 B-模，M 是一个 B-模。线性形式 \(u: L \to B\) 的给定等价于 B 中元素的一族 \(x = (x_i)_{i \in 1}\) ，对应关系为 \(x_i = u(e_i)\) 。作为 \(k\) -模复形， \(\mathbf{K}_{\mathrm{B}}^{\bullet}(u, \mathrm{M})\) （相应地， \(\mathbf{K}^{\mathrm{B}}(u, \mathrm{M})\) ）通过同构 \(\theta^{\bullet}\) （相应地， \(\theta\) .）与科斯居尔复形 \(\mathbf{K}'(x, \mathrm{M})\) （相应地， \(\mathbf{K}^{\bullet}(x, \mathrm{M})\) ）等同。例如 \(\mathbf{K}^{\mathrm{B}}(u)\) 与 \(\mathbf{K}'(x, \mathrm{B})\) 等同。

我们如同第 1 节（X，第 147 页）那样引入记号 H.(x, M)、H.(x, M) 等，由于模 A \(^{I}\) 是自由模，第 1 节和第 2 节的所有结果均可经必要修改后适用。例如，我们有同构

\[
\begin{array}{l} \mathrm{H} _ {0} (\boldsymbol {x}, \mathrm{M}) \to \mathrm{M} / (\boldsymbol {x}) \mathrm{M} \\ \mathrm{H} ^ {0} (\boldsymbol {x}, \mathrm{M}) \to \mathrm{Hom} _ {\mathrm{A}} \left(\mathrm{A} / (\boldsymbol {x}), \mathrm{M}\right), \end{array}
\]

其中 \((\pmb{x})\) 表示 A 的理想 \(\sum \mathbf{A}x_{i}\) 。又例如，若 \(\mathbf{K}.(\pmb{x},\mathbf{A})\) 在次数 \(>0\) 时零调，我们有同构

\[
\begin{array}{l}\mathrm{H} _ {r} (\boldsymbol {x}, \mathrm{M}) \rightarrow \operatorname{Tor} _ {r} ^ {\mathrm{A}} (k, \mathrm{M}),\\\operatorname{Ext} _ {\mathrm{A}} ^ {r} (k, \mathrm{M}) \rightarrow \mathrm{H} ^ {r} (\boldsymbol {x}, \mathrm{M}).\end{array}
\]

最后，假设 I 是（有限的且）全序的，例如 \(I = \{1, \ldots, n\}\) ；让我们把 \(\symbf{\Lambda}^{n}(\mathbf{A}^{\mathrm{I}})\) 借助基元素 \(e_{1} \wedge \ldots \wedge e_{n}\) 等同于 A，并翻译 X，第 149 页的同构 \(\mathbf{K}_{p}^{\mathrm{A}}(u, \mathrm{M}) \to \mathbf{K}_{\mathrm{A}}^{n-p}(u, \mathrm{M})\) 。通过 \((\theta_{p})\) 与 \((\theta^{n-p})\) 的搬运，它变为同构

\[
\mathrm{C} _ {\mathrm{I}} ^ {p} (\mathbf {M}) \rightarrow \mathrm{C} _ {\mathrm{I}} ^ {n - p} (\mathbf {M})
\]

它把 \(m \in C_{I}^{p}(M)\) 对应到满足下列条件的元素 \(\check{m}\) （属于 \(C_{1}^{n-p}(M)\) ）：

\[
\boldsymbol {m} \left(\alpha_ {1}, \dots , \alpha_ {p}\right) = \check {\boldsymbol {m}} \left(\beta_ {1}, \dots , \beta_ {n - p}\right)\tag{13}
\]

若 \((\alpha_{1},\ldots,\alpha_{p},\beta_{1},\ldots,\beta_{n-p})\) 是 \(\{1,\ldots,n\}\) 的一个偶置换。我们也注意到，当 \(I=\{1,\ldots,n\}\) 时，可以把 \(C_{I}^{p}(M)\) 与元素取自 M 的族 \(m(\alpha_{1},\ldots,\alpha_{p})\) （其中 \(\alpha_{1}<\alpha_{2}<\ldots<\alpha_{p}\) ）的集合等同；公式 (12) 仍然成立，关系式 (13) 亦然。

例。--- 取 \(\mathbf{M} = k[\mathbf{T}_1, \ldots, \mathbf{T}_n]\) ；科斯居尔复形 \(\mathbf{K}\) （ \(\partial/\partial\mathbf{T}\) , \(\mathbf{M}\) ）——它与自同态序列（ \(\partial/\partial\mathbf{T}_1, \ldots, \partial/\partial\mathbf{T}_n\) ）相伴——等同于 \(k[x_{1},\ldots,x_{n}]\) 在 \(k\) 上的德拉姆复形（X，第 44 页）：对 \(\boldsymbol{m}\in\mathbf{C}_{\{1,\ldots,n\}}^{p}(\mathbf{M})\) ，对应以微分形式

\[
\omega (\boldsymbol {m}) = \sum_ {\alpha_ {1} <   \dots <   \alpha_ {p}} \boldsymbol {m} \left(\alpha_ {1}, \dots , \alpha_ {p}\right) d x _ {\alpha_ {1}} \wedge \dots \wedge d x _ {\alpha_ {p}},
\]

参见公式 (12) 及第 44 页例 1。

